% !TEX root = ../main.tex
% !TEX program = lualatex
\documentclass[../main.tex]{subfiles}
\IfSubfilesClassLoaded{\externaldocument{\subfix{../build/main}}}{}
% ====================
% File: 17A_Research_and_Development.tex
% ====================
\begin{document}
\IfSubfilesClassLoaded{\chapter{Research \& Development}}{}
% === Research & Development (EDO 3.3.4) ===
\section{Research \& Development}

\subsection{Summary}
\begin{description}
  \item[\textbf{Technology posture.}] The U.S. relies heavily on industry research and commercial technologies while Navy labs retain a strong \ac{sandt} role in acquisition decisions~\autocite{edo-3-3-4-research-development-2025}.
  \item[\textbf{S\&T role.}] \ac{sandt} reduces integration risk by maturing technologies, conducting assessments, and advising requirements and acquisition communities~\autocite{edo-3-3-4-research-development-2025}.
  \item[\textbf{Continuum and readiness.}] The \ac{sandt} continuum runs from basic research to advanced technology; \ac{trl} maturity is central to insertion timing and risk~\autocite{edo-3-3-4-research-development-2025,USDRE-TRA-2025}.
  \item[\textbf{Transition mechanisms.}] \ac{atd} and \ac{jctd} accelerate technology transition with distinct objectives, oversight, and user involvement~\autocite{edo-3-3-4-research-development-2025}.
  \item[\textbf{Support partners.}] \acp{ffrdc} and \acp{uarc} are separate long-term support constructs that provide objective research, engineering, analysis, and systems acquisition support when the Department needs trusted technical depth outside the organic workforce~\autocite{edo-3-3-4-research-development-2025,DoW-FFRDC-UARC-2026}.
\end{description}

\subsection{R\&D Perspective and Posture}
\begin{description}
  \item[\textbf{Near term (today's Navy).}] Focused on current readiness and fleet needs led by the \ac{cno}, \ac{cmc}, \acp{ccmd}, Fleet, and the \ac{usmc}~\autocite{edo-3-3-4-research-development-2025}.
  \item[\textbf{Mid term (next Navy).}] Programs and resources shaped by \acp{peo}, \acp{syscom}, \ac{nro}, and \ac{opnav}~\autocite{edo-3-3-4-research-development-2025}.
  \item[\textbf{Long term (Navy the Nation needs).}] \ac{don} \ac{sandt}, the Naval War College, and research centers such as \ac{niwc}, \ac{nswc}, the Naval Air Warfare Center, \ac{nuwc}, \ac{nrl}, \ac{onr}, \ac{diu}, and \ac{darpa}~\autocite{edo-3-3-4-research-development-2025}.
\end{description}

\subsection{Using Commercial Technology}
\begin{description}
  \item[\textbf{Industry leverage.}] \ac{dod} \ac{sandt} budgets drive greater reliance on commercial investments while Navy labs retain an advisory role~\autocite{edo-3-3-4-research-development-2025}.
  \item[\textbf{Common pathways.}] \ac{osa}, \ac{duap}, \ac{ird}, \ac{mantech}, \ac{t2}, and \ac{sandt} affordability initiatives bring commercial technologies into defense systems~\autocite{edo-3-3-4-research-development-2025}.
\end{description}

\subsection{R\&D in the Acquisition Framework}
\begin{description}
  \item[\textbf{Technology opportunities.}] State-of-the-art assessments, experiments, and demonstrations inform program decisions~\autocite{edo-3-3-4-research-development-2025}.
  \item[\textbf{Insertion timing.}] New technology can enter in any phase, but the core maturity push is before Milestone~B: \ac{tmrr} reduces technology and integration risk, and official \ac{rdte} budget-activity guidance places advanced technology development and advanced component development/prototypes before Milestone~B~\autocite{edo-3-3-4-research-development-2025,DoDI5000-85,DoDFMR-Vol2B-Ch5,USDRE-TRA-2025}.
\end{description}

\subsection{Role of Science \& Technology}
\begin{description}
  \item[\textbf{Component S\&T executives.}] Evaluate battlefield deficiencies, establish \ac{sandt} projects, mature technologies, conduct independent assessments, and advise requirements and acquisition communities~\autocite{edo-3-3-4-research-development-2025}.
  \item[\textbf{Transition mechanisms.}] \ac{atd}, \ac{jctd}, experiments, and projects with explicit transition requirements move technology toward acquisition~\autocite{edo-3-3-4-research-development-2025}.
\end{description}

\subsection{FFRDCs and UARCs}
\begin{description}
  \item[\textbf{FFRDCs.}] Independent, not-for-profit, private-sector organizations established and funded to meet special long-term engineering, research, development, or analytic needs that cannot be met as effectively by Government or other private-sector resources. DoW-supported examples include CNA, IDA, RAND Arroyo Center, NDRI, Project AIR FORCE, Aerospace, MITRE NSEC, IDA Communications and Computing, MIT Lincoln Laboratory, and SEI~\autocite{edo-3-3-4-research-development-2025,DoW-FFRDC-UARC-2026}.
  \item[\textbf{UARCs.}] Nonprofit university-affiliated research organizations, not a subtype of \acp{ffrdc}; they maintain core engineering, research, and development capabilities tailored to long-term Department needs and are established, transferred, or terminated only by USW(R\&E). DoW-supported examples include Georgia Tech Research Institute, USC Institute for Creative Technologies, JHU Applied Physics Laboratory, Penn State Applied Research Laboratory, MIT Institute for Soldier Nanotechnologies, and the Systems Engineering Research Center~\autocite{edo-3-3-4-research-development-2025,DoW-FFRDC-UARC-2026}.
\end{description}

\subsection{Technology Readiness Levels}
\begin{itemize}
  \item TRL~1: Basic principles observed and reported~\autocite{edo-3-3-4-research-development-2025,USDRE-TRA-2025}.
  \item TRL~2: Technology concept or application formulated~\autocite{edo-3-3-4-research-development-2025,USDRE-TRA-2025}.
  \item TRL~3: Analytical and experimental proof of concept~\autocite{edo-3-3-4-research-development-2025,USDRE-TRA-2025}.
  \item TRL~4: Component or breadboard validation in laboratory environment~\autocite{edo-3-3-4-research-development-2025,USDRE-TRA-2025}.
  \item TRL~5: Component or breadboard validation in relevant environment~\autocite{edo-3-3-4-research-development-2025,USDRE-TRA-2025}.
  \item TRL~6: System/subsystem model or prototype demonstrated in relevant environment (desired for Milestone~B)~\autocite{edo-3-3-4-research-development-2025,USDRE-TRA-2025}.
  \item TRL~7: System prototype demonstrated in operational environment (desired for Milestone~C)~\autocite{edo-3-3-4-research-development-2025,USDRE-TRA-2025}.
  \item TRL~8: Actual system completed and qualified through test and demonstration~\autocite{edo-3-3-4-research-development-2025,USDRE-TRA-2025}.
  \item TRL~9: Actual system proven through successful mission operations~\autocite{edo-3-3-4-research-development-2025,USDRE-TRA-2025}.
\end{itemize}

\subsection{RDT\&E Budget Activities (6.1--6.8)}
\begin{description}
  \item[\textbf{6.1 Basic research.}] Foundational scientific study and experimentation without specific application to processes or products in mind~\autocite{edo-3-3-4-research-development-2025,DoDFMR-Vol2B-Ch5}.
  \item[\textbf{6.2 Applied research.}] Systematic expansion and application of knowledge to meet recognized military needs and determine feasibility of potential solutions~\autocite{edo-3-3-4-research-development-2025,DoDFMR-Vol2B-Ch5}.
  \item[\textbf{6.3 Advanced technology development.}] Component, subsystem, model, and technology-demonstration work to prove feasibility, operability, producibility, military utility, or cost-reduction potential before system-specific development~\autocite{edo-3-3-4-research-development-2025,DoDFMR-Vol2B-Ch5}.
  \item[\textbf{6.4 Advanced component development and prototypes.}] Integrated technology and prototype-system evaluation in realistic environments before Milestone~B, with TRL~6/7 maturity expected for major programs~\autocite{edo-3-3-4-research-development-2025,DoDFMR-Vol2B-Ch5}.
  \item[\textbf{6.5 System development and demonstration.}] Post-Milestone~B engineering and manufacturing development, integration, demonstration, and test work supporting Milestone~C and production decisions~\autocite{edo-3-3-4-research-development-2025,DoDFMR-Vol2B-Ch5}.
  \item[\textbf{6.6 RDT\&E management support.}] Test ranges, labs, test aircraft/ships, studies, analyses, and infrastructure or operations that support general \ac{rdte} efforts~\autocite{edo-3-3-4-research-development-2025,DoDFMR-Vol2B-Ch5}.
  \item[\textbf{6.7 Operational system development.}] Development upgrades for systems already fielded or approved for full-rate production, including projects with \ac{lrip} approval and expected production funding~\autocite{edo-3-3-4-research-development-2025,DoDFMR-Vol2B-Ch5}.
  \item[\textbf{6.8 Software and digital technology pilot programs.}] Software, electronic tools, systems, applications, services, requirements development, technical evaluations, and related activities directly supporting approved software and digital technology pilot programs~\autocite{DoDFMR-Vol2B-Ch5}.
\end{description}

\subsection{IRL and SRL}
\begin{description}
  \item[\textbf{Integration maturity.}] \ac{irl} and \ac{srl} complement \ac{trl} by measuring how well technologies integrate across interfaces~\autocite{edo-3-3-4-research-development-2025,USDRE-TRA-2025}.
  \item[\textbf{Nine-level scale.}] Like \ac{trl}, \ac{irl}/\ac{srl} use a 1--9 scale to communicate integration readiness and mission support maturity~\autocite{edo-3-3-4-research-development-2025,USDRE-TRA-2025}.
\end{description}

\subsection{Advanced Technology Demonstrations (ATDs)}
\begin{description}
  \item[\textbf{Purpose.}] Hardware/software prototypes that refine basic or applied research and demonstrate feasibility at relatively low cost~\autocite{edo-3-3-4-research-development-2025}.
  \item[\textbf{Funding and timing.}] Often \ac{onr}-funded (6.3) and usually pre-Milestone~B~\autocite{edo-3-3-4-research-development-2025,DoDFMR-Vol2B-Ch5}.
  \item[\textbf{User involvement.}] Limited or no user involvement; no leave-behind capability objective~\autocite{edo-3-3-4-research-development-2025}.
\end{description}

\subsection{Joint Concept Technology Demonstrations (JCTDs)}
\begin{description}
  \item[\textbf{Purpose.}] Evaluate military utility, develop \ac{conops} and doctrine, and establish realistic requirements~\autocite{edo-3-3-4-research-development-2025}.
  \item[\textbf{MUA.}] Produces a \ac{mua}; can substitute for an \ac{icd} for Milestone~A~\autocite{edo-3-3-4-research-development-2025}.
  \item[\textbf{Transition focus.}] Aims for early transition to acquisition or sustainment; may leave a go-to-war interim capability when appropriate~\autocite{edo-3-3-4-research-development-2025}.
\end{description}

\subsection{Small Business Innovation Research}
\begin{description}
  \item[\textbf{Purpose.}] \ac{sbir} stimulates innovation, meets federal R\&D needs, expands small-business participation, and commercializes federally funded innovations; \ac{sttr} adds the small-business/research-institution partnership lane~\autocite{edo-3-3-4-research-development-2025,USC-638-SBIR,SBIR-Apply-2026}.
  \item[\textbf{Funding cue.}] General \ac{sbir}/\ac{sttr} allocation rules are statutory and agency-specific. Current 15~U.S.C.~\S~638 requires 3.2\% for \ac{sbir} and generally 0.45\% for \ac{sttr}, while the FY~2025--FY~2031 \ac{dod} budget-calculation pilot uses at least 3.25\% \ac{sbir} and 0.46\% \ac{sttr} of the two-year average \ac{rdte} extramural budget~\autocite{USC-638-SBIR}.
  \item[\textbf{Phase I.}] Technical merit and feasibility; current SBIR.gov general guidance lists 6--12 months and \$50K--\$275K, with exact limits controlled by the solicitation and agency~\autocite{edo-3-3-4-research-development-2025,SBIR-Apply-2026}.
  \item[\textbf{Phase II.}] Continued research and development; current SBIR.gov general guidance lists about 24 months and \$400K--\$1.8M, with exact limits controlled by the solicitation and agency~\autocite{edo-3-3-4-research-development-2025,SBIR-Apply-2026}.
  \item[\textbf{Phase III.}] Commercialization or transition with no \ac{sbir}/\ac{sttr} funding~\autocite{edo-3-3-4-research-development-2025,SBIR-Apply-2026}.
\end{description}

\subsection{Quick Review}
\begin{itemize}
  \item What is the U.S. technology posture in defense acquisition~\autocite{edo-3-3-4-research-development-2025}? \textbf{Answer:} Heavy reliance on industry research and commercial technology.
  \item What is the role of \ac{sandt} in defense acquisition~\autocite{edo-3-3-4-research-development-2025,USDRE-TRA-2025}? \textbf{Answer:} Mature technology and reduce integration risk through research and assessment.
  \item When are new technologies most likely introduced in the life cycle~\autocite{edo-3-3-4-research-development-2025,DoDFMR-Vol2B-Ch5,USDRE-TRA-2025}? \textbf{Answer:} Prior to Milestone~B.
  \item What are key characteristics of \ac{atd}~\autocite{edo-3-3-4-research-development-2025}? \textbf{Answer:} Feasibility/maturity demonstrations and risk reduction.
  \item What are key characteristics of \ac{jctd}~\autocite{edo-3-3-4-research-development-2025}? \textbf{Answer:} Military utility evaluation, \ac{conops}/doctrine development, and potential go-to-war capability.
\end{itemize}

%====================
% End of file
%====================
\IfSubfilesClassLoaded{
  \RenewDocumentCommand{\entryname}{}{\textbf{\color{Modern} Acronym}}
  \RenewDocumentCommand{\descriptionname}{}{\textbf{\color{Modern} Definition}}
  \printnoidxglossary[
    type=\acronymtype,
    title=Acronyms,
    style=long-booktabs
  ]}{}
\end{document}
